\documentclass[12pt]{article}
\usepackage{amsmath}
\usepackage{amssymb}
\usepackage[margin=1in]{geometry}

\begin{document}

\begin{center}
    \textbf{\Large Homework}\\
    Tashfeen Hatim Said Omran\\
    UCF ID: 5862894\\
    Class: STA4321\\
    October 13, 2026
\end{center}

\vspace{0.5cm}

\noindent \textbf{Problem 3.92} \\
I define $Y$ as the number of trials until the first nondefective engine is found. $Y$ follows a geometric distribution with probability of success $p = 1 - 0.10 = 0.90$. The probability of failure (finding a defective engine) is $q = 0.10$.
I calculate the probability that the first nondefective engine is found on the second trial ($Y = 2$):
\[ P(Y = 2) = q^{2-1} p = (0.10)^1(0.90) = 0.09 \]

\vspace{0.5cm}

\noindent \textbf{Problem 3.94} \\
a) For the first nondefective engine, $Y \sim \text{Geometric}(p = 0.9)$. I find the mean and variance:
\[ E(Y) = \frac{1}{p} = \frac{1}{0.9} = \frac{10}{9} \approx 1.11 \]
\[ V(Y) = \frac{1-p}{p^2} = \frac{0.1}{(0.9)^2} = \frac{0.1}{0.81} = \frac{10}{81} \approx 0.123 \]
b) I define $W$ as the trial on which the third nondefective engine is found. $W$ follows a negative binomial distribution with $r = 3$ and $p = 0.9$. I find the mean and variance:
\[ E(W) = \frac{r}{p} = \frac{3}{0.9} = \frac{30}{9} = \frac{10}{3} \approx 3.33 \]
\[ V(W) = \frac{r(1-p)}{p^2} = \frac{3(0.1)}{(0.9)^2} = \frac{0.3}{0.81} = \frac{30}{81} = \frac{10}{27} \approx 0.370 \]

\vspace{0.5cm}

\noindent \textbf{Problem 3.110} \\
I define $Y$ as the number of nonlocal firms selected. $Y$ follows a hypergeometric distribution with population size $N = 6$, $r = 2$ nonlocal firms, and sample size $n = 3$. The total combinations are $\binom{6}{3} = 20$.
\[ P(Y = y) = \frac{\binom{2}{y}\binom{4}{3-y}}{\binom{6}{3}} \]
a) $P(Y = 1) = \frac{\binom{2}{1}\binom{4}{2}}{20} = \frac{2(6)}{20} = \frac{12}{20} = \frac{3}{5} = 0.6$ \\
b) $P(Y \ge 1) = 1 - P(Y = 0) = 1 - \frac{\binom{2}{0}\binom{4}{3}}{20} = 1 - \frac{1(4)}{20} = 1 - \frac{1}{5} = \frac{4}{5} = 0.8$ \\
c) $P(Y \le 1) = P(Y = 0) + P(Y = 1) = \frac{1}{5} + \frac{3}{5} = \frac{4}{5} = 0.8$

\vspace{0.5cm}

\noindent \textbf{Problem 3.122} \\
I define $Y$ as the number of customers arriving in an hour. $Y \sim \text{Poisson}(\lambda = 7)$.
\[ P(Y = y) = \frac{7^y e^{-7}}{y!} \]
a) $P(Y \le 3) = P(Y=0) + P(Y=1) + P(Y=2) + P(Y=3)$
\[ P(Y \le 3) = e^{-7} \left( \frac{7^0}{0!} + \frac{7^1}{1!} + \frac{7^2}{2!} + \frac{7^3}{3!} \right) = e^{-7} \left( 1 + 7 + 24.5 + 57.167 \right) = 89.667 e^{-7} \approx 0.0818 \]
b) $P(Y \ge 2) = 1 - P(Y \le 1) = 1 - e^{-7} \left( \frac{7^0}{0!} + \frac{7^1}{1!} \right) = 1 - 8e^{-7} \approx 0.9927$ \\
c) $P(Y = 5) = \frac{7^5 e^{-7}}{5!} = \frac{16807 e^{-7}}{120} \approx 0.1277$

\vspace{0.5cm}

\noindent \textbf{Problem 3.147} \\
I use the definition of the moment-generating function $m(t) = E(e^{tY})$. For a geometric distribution, $P(Y=y) = q^{y-1}p$.
\[ m(t) = \sum_{y=1}^{\infty} e^{ty} q^{y-1} p = \frac{p}{q} \sum_{y=1}^{\infty} (q e^t)^y \]
This is a geometric series with ratio $qe^t$. Assuming $qe^t < 1$, the sum evaluates to:
\[ m(t) = \frac{p}{q} \left( \frac{q e^t}{1 - q e^t} \right) = \frac{p e^t}{1 - q e^t} \]

\vspace{0.5cm}

\noindent \textbf{Problem 3.148} \\
I differentiate $m(t)$ to find $E(Y)$:
\[ m'(t) = \frac{d}{dt} \left[ p e^t (1 - q e^t)^{-1} \right] = p e^t (1 - q e^t)^{-1} + p e^t (-1)(1 - q e^t)^{-2}(-q e^t) = \frac{p e^t}{(1 - q e^t)^2} \]
\[ E(Y) = m'(0) = \frac{p e^0}{(1 - q e^0)^2} = \frac{p}{(1 - q)^2} = \frac{p}{p^2} = \frac{1}{p} \]
I differentiate $m'(t)$ to find $E(Y^2)$:
\[ m''(t) = \frac{d}{dt} \left[ p e^t (1 - q e^t)^{-2} \right] = p e^t (1 - q e^t)^{-2} + p e^t (-2)(1 - q e^t)^{-3}(-q e^t) = \frac{p e^t}{(1 - q e^t)^2} + \frac{2 p q e^{2t}}{(1 - q e^t)^3} \]
\[ E(Y^2) = m''(0) = \frac{p}{(1-q)^2} + \frac{2 p q}{(1-q)^3} = \frac{p}{p^2} + \frac{2 p q}{p^3} = \frac{1}{p} + \frac{2q}{p^2} = \frac{p + 2q}{p^2} = \frac{(1-q) + 2q}{p^2} = \frac{1+q}{p^2} \]
I find the variance:
\[ V(Y) = E(Y^2) - [E(Y)]^2 = \frac{1+q}{p^2} - \left(\frac{1}{p}\right)^2 = \frac{1+q-1}{p^2} = \frac{q}{p^2} \]

\vspace{0.5cm}

\noindent \textbf{Problem 3.153} \\
I match the given moment-generating functions to known distributions.\\
a) $m(t) = [(1/3)e^t + (2/3)]^5$. This matches the binomial MGF $(pe^t + q)^n$ with $n = 5$ and $p = 1/3$. The distribution is Binomial($n=5, p=1/3$). \\
b) $m(t) = \frac{e^t}{2 - e^t} = \frac{0.5 e^t}{1 - 0.5 e^t}$. This matches the geometric MGF $\frac{pe^t}{1 - qe^t}$ with $p = 1/2$. The distribution is Geometric($p=1/2$). \\
c) $m(t) = e^{2(e^t - 1)}$. This matches the Poisson MGF $e^{\lambda(e^t - 1)}$ with $\lambda = 2$. The distribution is Poisson($\lambda=2$).

\vspace{0.5cm}

\noindent \textbf{Problem 3.171} \\
For a Poisson distribution with mean $\mu = 100$, the variance is $\sigma^2 = \lambda = 100$, giving a standard deviation of $\sigma = 10$. \\
I use Chebyshev's theorem, setting $1 - 1/k^2 = 5/9$ to find the number of standard deviations $k$:
\[ \frac{1}{k^2} = 1 - \frac{5}{9} = \frac{4}{9} \implies k^2 = \frac{9}{4} \implies k = 1.5 \]
I calculate the interval $\mu \pm k\sigma$:
\[ 100 \pm 1.5(10) = 100 \pm 15 \]
The interval that includes at least $5/9$ of the sample values is $[85, 115]$.

\end{document}